% Lecture visuelle directe avec la vision du modèle ; aucun OCR.
% Source : PDF original, vue 027, côté gauche, page manuscrite 54.
% Statut : transcrit ; relecture indépendante à effectuer.
\chapter{Premiers rapports entre les futurs congréganistes. Rétablissement de la religion à Lyon}

\section*{Arrivée du cardinal Caprara}
% Mention marginale de la page manuscrite 54.
Nous vivions, ainsi que vous avez pu le remarquer dans le chapitre précédent, dans l'attente d'une délivrance prochaine, lorsque tout à coup on nous annonce l'arrivée du cardinal Caprara. Il étoit envoyé par le St Siège en qualité de légat à latere auprès du gouvernement français pour prendre de concert avec lui toutes les mesures nécessaires pour la publication et la mise à exécution du concordat.

Vous jugez quel effet dut produire dans une ville où aucun signe extérieur de religion n'étoit encore toléré et où les prêtres ne paroissoient point en public avec l'habit de leur état, l'apparition subite et inattendue d'un prince de l'Église revêtu de la pourpre romaine. Les autorités avoient reçu l'ordre de le recevoir avec tous les honneurs dus à sa haute dignité ; elles furent à sa rencontre et firent tirer le canon. Son Éminence séjourna à Lyon 24 heures ; toute la population accourut à l'hôtel de l'Europe où le cardinal étoit logé ; les salles, les galeries, la cour même de l'hôtel furent constamment remplies de nombreux chrétiens avides de recevoir la bénédiction d'un délégué du Saint-Siège, faveur d'autant plus appréciée qu'elle étoit alors plus extraordinaire.

Un légat à latere\footnote{Locution latine : « du côté ». Sens exact du titre à préciser, Q021. (Note du transcripteur).} en France, étoit pour l'opinion
% Source : PDF original, vue 027, côté droit, page manuscrite 55.
publique un fait d'une immense portée. Cet événement nous apportoit enfin la certitude complète de voir bientôt nos espérances se réaliser. Comme tout cela remuoit fortement les cœurs chrétiens ! Aussi est-ce à cette époque qu'on peut rapporter les premières ouvertures qui ont eu lieu entre les jeunes gens dévoués à la cause de la foi pour l'établissement d'une Congrégation. Depuis longtemps néanmoins le feu sacré couvoit sous les cendres. Nous avions été pendant de longues et douloureuses années les tristes témoins de l'astuce avec laquelle l'esprit de mensonge avoit organisé et discipliné des légions vraiment infernales dans l'exécrable but, hautement avoué, de détruire notre sainte religion, de renverser tous les trônes et de soumettre le genre humain au joug honteux et humiliant de l'anarchie. Cette triple conspiration savamment ourdie n'avoit que trop réussi ; nous avions vu de nos yeux le triomphe de l'impiété ; nous avions entendu les cris délirants de la société en proie à une licence effrénée ; le crime s'étoit appelé vertu ; la vertu s'étoit appelé crime. Un épouvantable déluge avoit tout bouleversé ; déluge de fausses doctrines qui avoit inondé les esprits, déluge de corruption qui avoit inondé les cœurs, déluge de sang au milieu duquel avoit presque disparu notre malheureuse patrie ; déluge de boue dans lequel elle s'étoit enfoncée... Alors, pauvres jeunes gens, échappés à ce naufrage presque universel par les soins d'une bienveillante Providence qui nous avoit placés dans
% Source : PDF original, vue 028, côté gauche, page manuscrite 56.
l'arche du salut sous l'égide assurée de la foi catholique, nous contemplions avec effroi l'horrible habileté avec laquelle l'esprit de ténèbres avoit dirigé cet odieux complot. Eh pourquoi, disions-nous alors, les mêmes moyens qui ont été employés pour détruire ne nous serviroient-ils pas à restaurer ? Ne pourrions-nous donc pas espérer avec le secours d'en haut de rétablir le règne de la vertu en agissant ainsi qu'on a fait pour établir celui du crime ? Ne seroit-ce pas que de faire servir à cette fin leurs propres armes ? Et Dieu aidant, ne pourrions-nous pas obtenir la victoire ? C'est ainsi que le Très-Haut avoit placé dans le cœur de plusieurs d'entre nous la sainte et salutaire pensée d'employer à procurer le règne du Seigneur et le triomphe de la foi les richesses d'iniquité amassées par les artisans du crime.

\section*{Premières ouvertures}
% Mention marginale de la page manuscrite 56.
Dès le moment où se manifestèrent les premières lueurs d'espérance de paix pour l'Église, M. Jordan, aujourd'hui curé de St Bonaventure, alors simple laïque, avoit essayé, sous la direction de M. de Villers, vicaire général du diocèse de Lyon, saint prêtre et zélé missionnaire, et concurremment avec celui qui tient ici la plume, de former un noyau de Congrégation. Mais hélas ! les moments marqués par la Providence n'étoient point encore arrivés. Tous leurs efforts échouèrent complètement. La timide colombe avoit essayé de sortir de l'arche, mais les eaux de
% Source : PDF original, vue 028, côté droit, page manuscrite 57.
l'iniquité n'étoient point suffisamment retirées. Elle ne put point trouver une terre assez sèche pour pouvoir s'y arrêter. Il fallut donc qu'elle rentrât et qu'elle attendît qu'un nouveau signal fût donné par le maître souverain. Nous allons la voir s'élancer une seconde fois ; mais cette fois-ci elle ne reparoîtra qu'avec le rameau d'olivier, c'est-à-dire avec cette alliance contractée avec le ciel sous les auspices de Marie qui nous est donnée comme le signe de la paix que Dieu veut bien accorder aux hommes de bonne volonté.

Cherchons donc maintenant, à suivre jusque dans leurs plus petits détails ces circonstances que leurs résultats rendent si intéressantes pour nous. Je voudrois connoître, pour vous le transmettre, tout ce qui s'est passé alors dans le cœur de la plupart des congréganistes de la première année ; je ne doute nullement que nous y retrouverions les mêmes sentimens qui animoient ceux de mes amis entre lesquels les premières paroles qui ont amené la fondation de la Congrégation ont été échangées. Mais dans l'ignorance où je suis d'une grande partie de ce qui s'est dit ou fait à cette époque, je ne puis que rassembler mes souvenirs pour vous les faire parcourir avec moi. D'autres lignes ont pu aboutir au même centre ; mais je suis certain que celle-là nous y conduira directement.

J'avois fait vers la fin de 1801 ou au commencement de 1802 connoissance avec un jeune homme qui s'appeloit M. Delagrange\footnote{Graphie du nom à confirmer sur les occurrences suivantes, Q023. (Note du transcripteur).}. Nous nous étions trouvés dans les oratoires. Il étoit fort pieux, ardent et zélé pour la cause de la religion.
% Source : PDF original, vue 029, côté gauche, page manuscrite 58.
Il éprouvoit avec beaucoup de vivacité tous les sentimens que les circonstances du moment faisoient naître dans les âmes. Une même manière de voir, de penser, de sentir devoit naturellement nous rapprocher. Aussi chaque dimanche à l'issue des offices nous cherchions à nous rejoindre et nous faisions souvent ensemble de longues promenades. Notre liaison avoit eu pour base ce désir commun que nous éprouvions de voir complètement triompher notre foi. Ces entretiens réitérés entre deux jeunes gens animés d'un même esprit eurent bientôt amené la confiance ; aussi nous en vînmes au point de nous communiquer nos pensées les plus intimes. Nos conversations rouloient presque toujours sur les circonstances du moment. Nous nous réjouissions bien sincèrement d'un retour à la religion qui ne paroissoit plus douteux, mais en même temps, plus que jamais, nous comprenions la nécessité de nous unir, de nous soutenir, de nous mettre en état de pouvoir à l'avenir défendre notre foi ; avec quelque efficacité, en un mot, de former entre nous une sainte ligue pour le bien, comme les méchants avoient formé une ligue infernale pour le mal. C’est par ces épanchements de l’amitié que le Seigneur préparoit de plus en plus nos cœurs à recevoir la grâce qu'il nous destinoit. Chaque promenade nous faisoit faire un pas en avant, si ce n'étoit pas encore d'une manière active, c'étoit du moins pour les dispositions où elle plaçoit nos cœurs. Après les deux ou trois premières conversations, M. Delagrange me parla de M. Ripoud et me pria de lui permettre de lui faire part
% Source : PDF original, vue 029, côté droit, page manuscrite 59.
de ce qui faisoit l'objet de nos conversations qu'on pouvoit déjà considérer comme parvenu à l'état de projet. J'étois moins à portée de voir habituellement M. Ripoud que M. Delagrange, mais je le connoissois parfaitement ; je savois tout son dévouement à la sainte cause de la religion ; je savois aussi que son zèle étoit prudent et éclairé ; ce fut donc avec plaisir que je consentis à le voir mettre en tiers dans nos secrets. M. Ripoud accueillit très bien cette première ouverture, mais vous le savez comme moi, Messieurs, il étoit calme et prudent ; pour opérer avec plus de sagesse, il ne vouloit rien précipiter. Il se mit au contraire en devoir d'arrêter l'impétuosité de M. Delagrange dont le zèle ardent ne connoissoit aucun obstacle et qui auroit désiré que sans autre examen on mît de suite la main à l'œuvre.

De mon côté, avant d'aller en avant, j'avois cru qu'il étoit nécessaire de m'appuyer des conseils d'hommes éclairés. J'eus la pensée de m'adresser à une personne respectable en qui j'avois toute confiance et que je considérois comme une autorité d'autant plus compétente pour me diriger qu'il avoit été pendant plusieurs années président du conseil diocésain. Sa voix pouvoit être considérée par nous comme celle de l'Église elle-même. Cet homme étoit M. Linsolas, le vicaire général dont nous avons déjà parlé qui avoit déployé tant d'habileté dans son administration. Une persécution particulière à sa personne venoit de nous l'enlever ; il étoit alors prisonnier sur parole à Turin ; je lui écrivis donc pour demander conseil. Ce
% Source : PDF original, vue 030, côté gauche, page manuscrite 60.
vénérable ecclésiastique nous encouragea beaucoup à donner suite à ce pieux projet, mais en même temps il nous conseilloit d'attendre pour mettre la main à l'œuvre le moment où le concordat seroit mis à exécution. Il supposoit que cette époque devoit nécessairement provoquer dans la ville un mouvement religieux qui ne pouvoit que favoriser beaucoup l'exécution de nos desseins. Je communiquai ce conseil à M. Ripoud ; il devint la règle de sa conduite et de la mienne. Le jeune M. Delagrange en étoit tout dolent ; mais malgré toute son ardeur il se vit obligé d'attendre. Cependant nous ne manquions pas de nous revoir chaque dimanche, et la Congrégation future étoit le plus souvent l'objet de nos entretiens ; nous donnions un libre cours à nos imaginations ; nous eussions voulu pouvoir pénétrer dans la profondeur de l'avenir pour y découvrir les destinées de cet établissement tout prêt à naître. Pauvres enfants que nous étions ! Foibles instruments de la miséricorde de notre Dieu ! Pouvions-nous prévoir alors cette immensité de grâces que nous devions voir par la suite descendre d'années en années sur nous et sur nos frères !

Ces rapports d'amitié n'existoient pas seulement entre nous, ils s'étendoient sur un grand nombre de jeunes gens qui tous ou presque tous sont devenus congréganistes. Ils formoient déjà entre eux, sans avoir contracté encore aucun engagement, comme une espèce de société dont le but étoit d'accélérer par tous les moyens
% Source : PDF original, vue 030, côté droit, page manuscrite 61.
qui dépendoient d'eux le retour de la religion dans notre ville. Déjà l'année précédente les uns et les autres s'étoient employés pour ouvrir à la piété des fidèles de grands oratoires. Cette année, devenus plus ambitieux, nous cherchions les moyens de rendre enfin au culte du Seigneur les églises que la foi de nos pères avoit élevées. La prise de possession d'une église étoit un pas que le triomphe de la religion faisoit en avant ; aussi tel étoit le but constant de nos désirs et de nos efforts.

\section*{Ouverture des Chartreux}
% Mention marginale de la page manuscrite 61.
Une des premières circonstances de ce genre et celle qui fut peut-être l’occasion [lecture incertaine : de la joie simultanée] du plus grand nombre de futurs congréganistes\footnote{Lecture du groupe après « l’occasion » à confirmer, Q025. (Note du transcripteur).} fut l'ouverture de l'église des Chartreux. La Providence qui destinoit cette sainte maison à être si souvent le théâtre des faveurs qu'elle réservoit à la Congrégation vouloit d'avance nous en faire prendre possession. La nouvelle se répand parmi les fidèles que par suite d'arrangements pris avec le propriétaire l'église des Chartreux a été rendue à la religion. Aussitôt les amis (je crois pouvoir déjà désigner par ce nom les futurs congréganistes), les amis donc s'empressent d'accourir, balayent, nettoyent, ornent l'église, tout cela leur paroît des fonctions honorables. C'est en même temps pour eux une véritable fête. Bientôt tout est disposé. On jette un coup d'œil pour voir s'il ne manque rien. Mais, dit l'un de nous, il n'y a point de chaire ; j'en ai aperçu une dans l'intérieur de la maison, répond un autre. À ces mots tous les amis s'élancent, enlèvent la chaire, la chargent sur nos épaules, la transportent dans l'église et la fixent à sa place. Tout cela
% Source : PDF original, vue 031, côté gauche, page manuscrite 62.
est l'affaire d'un instant. Cette chaire alloit devenir la chaire de vérité ; c’est du haut de cette tribune sacrée que la bonne nouvelle alloit de nouveau être annoncée au peuple chrétien, il nous sembloit donc en la replaçant [mot illisible] la religion elle-même dans notre ville\footnote{Raccord et verbe de cette proposition incertains, Q026. (Note du transcripteur).}. Les amis assistèrent en grand nombre, entourés d'un concours nombreux de fidèles, à cette réconciliation d'église, prélude de tant d'autres.

Les amis de la religion étoient alors bien consolés ; une chose toutefois les affligeoit bien vivement encore : c'étoit la présence des prêtres schismatiques dans nos églises et l'obstination déplorable d'un grand nombre d'entre eux. Les chefs de ce parti surtout, que leur orgueil tenoit toujours éloignés de la voie droite, ne cessoient de répéter que ce ne seroit pas eux qui iroient au pape, mais que le pape viendroit à eux. Malgré cela, comme dans tous les temps et même au plus fort de la Terreur, il y avoit eu des rétractations et il y en avoit toujours. Elles devenoient même plus fréquentes car la crainte qui jusque-là avoit pu retenir les hommes timides étoit un motif qui ne subsistoit plus. Il y avoit cependant quelque mérite dans ces retours ; il falloit bien que ces prêtres repentants fussent guidés par les considérations de la foi, car en attendant la publication du concordat ils pouvoient se considérer comme certains d'obtenir leur réconciliation sans aucune espèce d'humiliation, tandis que le conseil diocésain, ferme sur les principes, ne les recevoit à pénitence que sous la condition de faire abjuration dans un oratoire en présence de
% Source : PDF original, vue 031, côté droit, page manuscrite 63.
l'assemblée des fidèles et d'être ensuite réduits à la communion laïque pendant un temps plus ou moins long. On abrégeoit en raison de la ferveur du pénitent la durée de cette espèce de pénitence canonique.

\section*{Rétractations. Jour de Pâques}
% Mentions marginales de la page manuscrite 63.
J'ai assisté à une de ces abjurations ; permettez-moi de vous en parler ici : c'étoit le P. Roger qui la recevoit. Cela eut lieu le jour de Pâques 1802. Quel beau jour, Messieurs, que ce jour de Pâques 1802 ! Comme autrefois il l'avoit fait entendre à Lazare dans le tombeau, J.-C. venoit encore faire entendre sa voix puissante à l'Église de France : Sortez de cet obscur tombeau où la persécution vous a renfermée ; [lecture incertaine : en liberté il paroisse] au grand jour\footnote{Lecture de l'injonction incertaine, Q027. (Note du transcripteur).} ; soyez désormais délivrée des liens qui vous tenoient captive ; levez-vous, marchez en assurance.

Voilà vraiment ce qui se passoit dans ce grand jour, car le légat du St Siège rouvroit solennellement, au nom du souverain pontife, la métropole de Paris, et on proclamoit le concordat dans toute la France. Aussi l'Église répétoit-elle avec transport les paroles de l'office de cette grande solennité : Voilà le jour que le Seigneur a fait, réjouissons-nous et tressaillons d'allégresse.

Eh bien ! ce jour-là sur les onze heures du matin je me trouvois dans un petit et pauvre oratoire chez un ouvrier en soie à la rue St Georges. Le P. Roger est à l'autel ; un prêtre en soutane est à genoux au milieu de l'assemblée ; il demande d'une voix entrecoupée de sanglots le pardon et la pénitence. Les fidèles vivement émus peuvent à peine contenir les sentimens dont ils sont agités. Le
% Source : PDF original, vue 032, côté gauche, page manuscrite 64.
P. Roger partage l'émotion générale ; il adresse des paroles de paix au schismatique repentant ; le félicite de la généreuse résolution qu'il a prise de rentrer enfin dans le sein de l'Église de J.-C., de cette Église toujours assise sur la base inébranlable de la vérité éternelle dont seule elle est dépositaire ; c’est cette immuable vérité qui dans tous les siècles a triomphé de tous les schismes et de toutes les hérésies, c’est cette immuable vérité, peuple fidèle qui va bientôt rouvrir les portes de vos temples et relever vos autels : \textit{et veritas Domini manet in æternum}\footnote{Traduction littérale : « Et la vérité du Seigneur demeure éternellement. » Référence de la citation à vérifier, Q028. (Note du transcripteur).}. Ces paroles ne sont pas achevées que le bourdon de notre cathédrale, muet depuis tant d'années, se fait entendre. Ses mâles vibrations font vibrer tous nos cœurs, car il annonce qu'en ce moment même le traité de paix et de réconciliation avec notre Dieu se publie solennellement sur nos places publiques. Faites-vous, si vous le pouvez, Messieurs, car je n'ai point d'expressions capables de le peindre, faites-vous quelque idée de ce qui se passoit alors et dans l'âme du ministre de paix chargé par l'Église de la donner à cet humble pénitent et dans l'âme du prêtre revenu de ses erreurs dont la bonne volonté a mérité de recevoir cette paix envoyée par le Ciel et enfin dans les âmes des fidèles assemblés autour de l'autel de J.-C. qui sont les héritiers et les successeurs des chrétiens des premiers siècles et sont en ce moment les témoins d'un spectacle semblable à ceux qu'on voyoit bien souvent dans les beaux jours de la primitive Église.

% Source : PDF original, vue 032, côté droit, page manuscrite 65.
Vous devez bien le penser, Messieurs, cette fête de Pâques déjà si belle par elle-même avoit pour nous cette année-là quelque chose de plus solennel qu'à l'ordinaire. Des prières d'actions de grâces avoient été ordonnées par les vicaires généraux. Un \textit{Te Deum}\footnote{Début de l'hymne latine : « Vous, Dieu, [nous vous louons] ». (Note du transcripteur).} devoit être chanté dans tous les oratoires. Vous jugez avec quel bonheur les amis, futurs congréganistes, mêloient les accents de leur reconnoissance aux voix du peuple fidèle. Je pourrois en citer qui, ne pouvant se lasser de témoigner leur joie, passoient d'un oratoire à l'autre et ont chanté le \textit{Te Deum} jusqu'à quatre fois.

\section*{Réconciliation de St Jean}
% Mention marginale de la page manuscrite 65.
Dans le travail préparé d'avance par les deux puissances pour la nomination des nouveaux évêques de France, Mgr Fesch, oncle du premier consul, avoit été désigné pour le siège de Lyon. Mais comme depuis bien des années il n'exerçoit plus les fonctions sacerdotales, il voulut se préparer à l'épiscopat par une retraite de six mois au séminaire de St Sulpice à Paris. On ne pouvoit pas laisser pendant un aussi long temps sans pasteur un diocèse de cette importance. Pour obvier à cet inconvénient on chargea Mgr de Mérinville, évêque de Chambéry, de son administration. Ce prélat arriva à Lyon quelques jours avant la fête de la Pentecôte et de concert avec lui on choisit l'époque de cette grande solennité pour rendre à la religion l'église primatiale de St Jean.

La veille à six heures du soir eut lieu la réconciliation de l'église. Vous jugez si les amis s'empressèrent d'y assister.
% Source : PDF original, vue 033, côté gauche, page manuscrite 66.
Avant la cérémonie, l'église étant encore considérée comme un lieu profane, les fidèles qui s'y trouvoient en grand nombre conservoient le chapeau sur la tête comme au milieu d'une place publique. Le moment arrive enfin où ce temple auguste devoit retrouver son antique sainteté. Un prêtre s'avance et d'une voix forte adresse ces trois paroles à la multitude assemblée : Silence, respect, attention. À l'instant toutes les conversations cessent, toutes les têtes sont découvertes. L'étendard de la croix suivi d'un nombreux clergé vient prendre au nom de J.-C. possession de l'église. Le bourdon de la métropole se fait entendre au même moment comme si le Seigneur eût voulu manifester son intervention directe à la réconciliation d'une basilique si féconde en souvenirs, un éclair répand une lumière éclatante dans la profondeur des voûtes ; la foudre vient mêler son imposante voix aux mâles vibrations de l'airain sacré et les eaux du ciel tombent par torrents ; on croiroit qu'elles ont reçu la mission de purifier l'extérieur du temple pendant que l'eau bénite mêlée aux prières de l'Église en purifie l'intérieur. Lorsque les prières de la réconciliation sont achevées, les jeunes amis qui dans cette intention s'étoient emparés d'avance de la croix et des chandeliers s'empressent de les élever sur l'autel. C'étoit toujours une petite part que nous prenions au rétablissement du culte du Seigneur. Il est à remarquer que dans la plupart des cérémonies
% Source : PDF original, vue 033, côté droit, page manuscrite 67.
de ce genre qui se sont succédées ce sont presque toujours nos amis qui ont rempli cette fonction et qui l'ont remplie les uns avec les autres. Ne pourroit-on pas dire que Dieu vouloit nous avertir par là que nous étions destinés à travailler à sa gloire en réunissant nos efforts.

La ville entière étoit sur pied pour assister à la cérémonie du lendemain. Mgr de Mérinville officia pontificalement dans la métropole qui venoit d'être rendue à la religion. Tout le clergé catholique de Lyon étoit réuni dans le chœur. Plus de trois cents prêtres étoient présents. Qu'il étoit touchant de voir reparaître en un jour de paix et de réconciliation ces vénérables confesseurs de la foi ! La plupart d'entre eux avoit généreusement supporté pour la défense de la religion tous les genres de privations. Plusieurs avoient échappé miraculeusement à la mort ; d'autres revenoient de Rochefort ou de Sinnamary où ils avoient enduré les traitements les plus barbares ; presque tous étoient encore revêtus de leur costume laïc car ils n'avoient pas eu le temps de se procurer de soutane ; un grand nombre n'avoit pas même coupé leurs cheveux et les portoient en queue suivant l'usage du temps. En contemplant cette sainte assemblée, spectacle touchant et vraiment digne des anges, les fidèles se sentoient émus et pénétrés d'admiration ; ils avoient devant les yeux l'Église des catacombes qui se montroit au grand jour glorieuse et triomphante après avoir bien combattu. Quelques
% Source : PDF original, vue 034, côté gauche, page manuscrite 68.
membres de ce digne clergé avoient été désignés pour conduire le chant ; ceux-là étoient seuls en surplis et occupoient des places à part. Quant aux laïques, parmi lesquels sans que nous le répétions vous comprenez bien, je pense, doit entrer la Congrégation future, ils avoient formé deux chœurs dans le bas de l'église. De là leurs voix répondoient aux voix du clergé ; c'étoit avec un saint enthousiasme que l'assemblée des fidèles toute entière répétoit à l'envi les louanges du Seigneur. Lorsqu'on entonna le chant majestueux du Credo, il nous sembla voir l'illustre Église de Lyon promulguer de nouveau la perpétuité de la foi par la bouche de ses enfants. Entendez leur profession de foi ; c'est toujours la foi des Pothin et des Irénée ; c'est toujours l'antique symbole que les Pères de Nicée et de Constantinople ont rédigé et que les fidèles nous ont transmis. Mais voyez, d'admirables paroles ! Comme elles expriment clairement et notre foi toujours immuable et notre attachement inébranlable à l'Église de J.-C. : \textit{et unam sanctam Catholicam et Apostolicam Ecclesiam}\footnote{Traduction littérale : « Et [je crois en] une seule Église sainte, catholique et apostolique. » (Note du transcripteur).}. Ces mots aussi dans le transport qui nous anime, nos voix ne peuvent plus suffire toutes seules ; nos cœurs, nos âmes, tout ce qui vit et respire en nous s’élancent à la fois pour les faire entendre ; les échos de la vieille basilique en sont émus ; ils repoussent avec indignation les accents [mot incertain : chimères]\footnote{Qualification des accents à confirmer, Q037. (Note du transcripteur).} que le schisme ou le paganisme leur ont fait naguère répéter et nous répondent par la même profession de foi. Il y a
% Source : PDF original, vue 034, côté droit, page manuscrite 69.
longtemps qu'ils la connoissent. L'Église entière assemblée dans le même lieu au treizième siècle la leur avoit fait entendre avec énergie au moment où les Grecs reconnurent, au moins un instant, l'unité catholique.

Les catholiques devenus un peu trop ardents dans l'ivresse de leur triomphe s'étoient permis à l'occasion de cette solennité une petite licence qui n'étoit peut-être pas tout à fait suivant les règles de la prudence. Ils avoient voulu sinon exclure les prêtres schismatiques qui ne s'étoient pas rétractés, du moins les réduire à la condition des pénitents des premiers siècles en les forçant à s'arrêter au bas de l'église. Des barrières avoient été placées et des laïques dévoués s'étoient chargés de les garder et d'empêcher tout prêtre constitutionnel de passer outre. Un obstiné trouva cependant je ne sais comment moyen de se soustraire à cette exclusion et il parvint au chœur. Il étoit en costume complet, soutane, surplis, rabat, rien n’y manquoit. Il va se mêler à la foule des prêtres ; ceux-ci le voyant en surplis l'envoient rejoindre les chantres ; il se présente donc ; mais ici le petit bout d'oreille échappé par malheur découvrit la fausseté et l'erreur ; parmi tous les prêtres catholiques présents on n'auroit pas pu alors en trouver un qui eût un seul rabat à sa disposition. À la vue de malencontreux rabat la position équivoque de celui qui le portoit est reconnue par les chantres ; ils le repoussent à leur tour. Ce petit incident pouvoit donner lieu à un véritable scandale lorsqu'un des laïques délégués par l'autorité pour veiller à l'ordre de la cérémonie, s'apercevant de ce qui se passoit, pria
% Source : PDF original, vue 035, côté gauche, page manuscrite 70.
ce pauvre prêtre de se rendre à la sacristie. Là, après lui avoir reproché son imprudence, il lui fit quitter son habit de chœur et l'engagea à se retirer. Cette scène eut un heureux résultat. Celui qui en avoit été l'objet éprouva une humiliation salutaire d'être vu ainsi expulsé de l'assemblée des fidèles ; la réflexion excita les remords ; ceux-ci produisirent une bonne résolution, et dès le soir même il s'empressa de préparer sa réconciliation avec l'Église en envoyant sa rétractation.

Vous vous imaginez sans doute, Messieurs, l'impression profonde que produisoit sur tous les cœurs chrétiens et en particulier sur les nôtres cette belle fête de l'ouverture de St Jean. Dieu nous tenoit encore en réserve de nouvelles douceurs. Elles ne manquoient pas pour ceux qui aimoient la religion dans ces jours de bénédiction. Nous ne nous trouvions alors jamais plus heureux que lorsqu'une solennité nouvelle nous rassembloit dans les temples du Seigneur. L'octave de la Fête-Dieu vint nous en procurer l'occasion. Les desservants de l'église des Chartreux favorisés par la position isolée de cette église pensèrent à en profiter pour faire une procession solennelle du Saint-Sacrement qu’il n’eût pas été possible alors de déployer dans l'intérieur de la ville. Le quartier des Chartreux n'étoit pas habité comme il l'est maintenant ; par ses alentours cette église ressembloit tout à fait à une église de campagne. Les cloîtres qui l'environnent existoient dans leur entier et offroient toutes les facilités désirables pour qu'une
% Source : PDF original, vue 035, côté droit, page manuscrite 71.
procession pût les parcourir sans effaroucher les amis du monde, [lecture incertaine : l’œil vigilant] d’une police inquisitoriale\footnote{Raccord après « amis du monde » incertain, Q034. (Note du transcripteur).}. Nous avions tous éprouvé une peine bien vive en nous voyant forcés de passer le jour solennel consacré par l'Église à la mémoire du mystère de l'amour de notre Sauveur pour nous, sans avoir pu lui rendre en la présence du peuple un hommage éclatant de notre foi et de notre reconnoissance. Nous trouvâmes bien doux de pouvoir nous dédommager de ce sacrifice en terminant l'octave. Une pieuse conspiration est donc formée. Tous les amis doivent se rendre aux Chartreux ; ils y arrivent munis d'un cierge auquel un bouquet de fleurs est attaché ; c'étoit le signe de ralliement convenu. Ah ! Messieurs, combien on est sensible à ces jouissances de l'âme lorsque on a été réduit à en supporter longtemps la privation ! De quelle joie pure n'étions-nous pas pénétrés lorsque nous vîmes s'ébranler la pompe triomphale dont nous-mêmes allions être appelés à faire partie ! Une cohorte de vierges chrétiennes vêtues de blanc réunies sous l'étendard de la Vierge sans tache ouvre la marche ; viennent ensuite de nombreux fidèles ; après eux les prêtres du Seigneur précédés par la croix, par cette croix qui a sauvé le monde, par cette croix qui vient de triompher encore comme elle a triomphé dans tous les siècles passés et comme elle triomphera dans tous les siècles à venir, par cette croix par laquelle le Christ règne dans le temps et régnera dans l'éternité. Au milieu de ce vénérable clergé
% Source : PDF original, vue 036, côté gauche, page manuscrite 72.
paroissent de pieux lévites revêtus d'une tunique blanche, symbole d'innocence et de paix. Ils balancent devant la majesté suprême l'urne des parfums et font élever en sa présence des nuages d'encens auxquels viennent se mêler les fleurs que de jeunes enfants sèment à pleines mains sur les pas du Dieu vivant. Le soleil de justice rayonnant de lumière, l'adorable Eucharistie franchit les portes de son temple ; elle est environnée par ces jeunes gens qui bientôt formeront la famille de Marie ; leurs flambeaux allumés sont un symbole de la foi qu'ils ont conservée et de la charité dont le Sauveur du monde veut que leurs cœurs soient embrasés ; ces fleurs dont ces flambeaux sont ornés figurent cette bonne odeur de J.-C. dont ils doivent se parfumer car bientôt ils vont être appelés à la répandre. C'est ainsi que guidés par notre bon maître dont ils forment l'escorte ces jeunes gens vont parcourir ces cloîtres habités jadis par de pieux cénobites et dans lesquels le bon pasteur réserve de délicieux pâturages à un petit troupeau composé de brebis privilégiées qu'il choisira lui-même. Ces voûtes autrefois témoins muets des pieux soupirs des enfants de St Bruno retentissent aujourd'hui des chants d'allégresse d'un peuple fidèle ; encore un peu de temps et elles vont entendre répéter avec enthousiasme le refrain chéri de la charité lorsqu'il exaltera la douce union qui va rassembler les enfants de Marie comme autant de frères : \textit{Ecce quam bonum et quam jucundum,
% Source : PDF original, vue 036, côté droit, page manuscrite 73.
habitare fratres in unum}\footnote{Traduction littérale : « Voici combien il est bon et agréable que des frères habitent ensemble. » Référence à vérifier, Q036. (Note du transcripteur).}.

Ce fut pour nous tous un jour de bonheur que le jour où il nous fut donné d'assister à cette procession ; un sentiment délicieux remplissoit nos âmes ; Jésus par sa douce présence nous combloit d'une joie ineffable ; le souvenir du passé et les jouissances du présent nous convioient également à célébrer les miséricordes du Seigneur. Quels transports n'aurions-nous pas éprouvés si nous eussions connu l'avenir ? Oh ! c'est bien alors que nous eussions pu nous écrier : Oui, Jésus-Christ qui étoit hier, qui sera demain et dans les siècles des siècles, ne cesse pas d'étendre d'âge en âge sa miséricorde sur ceux qui le craignent ; il ouvre constamment ses bras pour y recevoir ses disciples bien-aimés, et s'ils s'approchent de lui avec une confiance toute filiale, il les fait reposer sur son cœur et les embrase de son amour.

Toutes ces grâces répandues sur nous avec tant d'abondance nous disposoient de plus en plus à nous consacrer entièrement à travailler à la gloire de notre divin maître. Aussi nous approchons du moment où la Congrégation fidèle à la voix qui l'appelle va enfin entrer dans la carrière qui lui est réservée ; nous allons donc terminer ce précis historique pour commencer à nous occuper de l'histoire de cette société. Mais auparavant résumons ce que nous avons dit jusque ici afin d'embrasser d'un seul coup d'œil les moyens qu'il a plu à la Providence d'employer pour préparer les
% Source : PDF original, vue 037, côté gauche, page manuscrite 74.
premiers congréganistes à la mission qu'elle vouloit leur confier. Vous avez pu remarquer qu'en cherchant à développer les causes qui ont produit cet établissement nous avons été amenés à suivre les points principaux de l'histoire de notre Révolution. C'est qu'en effet c'est par le spectacle qu'ils ont offert que Dieu a voulu disposer les cœurs qu'il destinoit à accomplir ses desseins. Il les a rendus les témoins de l'horrible succès des impies coalisés contre la religion et la société ; il a excité en eux en déroulant à leurs yeux ce lamentable tableau, d'une part une salutaire indignation contre les doctrines qui ont occasionné tant de maux, de l'autre le noble désir d'y apporter un remède. Pour que l'iniquité triomphante ne parvînt pas à ébranler leur foi, il leur a fait voir en même temps dans tout son éclat l'accomplissement des promesses qui ont été faites à l'Église de J.-C. C'est par ces moyens que dans sa miséricorde le Seigneur daignoit inspirer à plusieurs d'entre ces jeunes chrétiens le désir bien sincère de se dévouer pour la défense de sa cause. Mais leur bonne volonté ainsi isolée n'auroit pas suffi au but que le Tout-Puissant se proposoit. Leur mission devoit être exercée collectivement ; Dieu leur a fait voir par l'exemple des ennemis de l'Église quelle est la force d'une action commune ; l'orgueil a pu être le mobile de ces sociétés perverses ; mais la charité seule pouvoit être le lien de la société
% Source : PDF original, vue 037, côté droit, page manuscrite 75.
des enfants de Dieu ; elle ne pouvoit être fondée que par la charité et maintenue que par la charité, et pour donner plus d'élan à son action cette charité doit prendre ici le caractère d'une amitié chrétienne, d'une véritable fraternité.

Après avoir rempli les cœurs de zèle pour sa gloire, Dieu a donc voulu y faire naître aussi cet autre sentiment. Il a voulu que ce désir de ne faire qu'un avec ses frères régnât en maître dans ces âmes qu'il préparoit pour accomplir son œuvre. Pour augmenter la vivacité de ce désir il les a tenus dans l'état d'isolement où la persécution les avoit placés ; il leur a fait ainsi regretter les charmes de cette amitié chrétienne ; pour rendre plus vive encore la soif qu'ils en éprouvoient il a fait distiller d'abord jusqu'à eux quelques gouttes de cette délicieuse liqueur ; il leur a permis seulement d'y tremper leurs lèvres ; après les avoir ainsi accoutumés à en savourer les douceurs il l'a enfin laissé couler par torrents au milieu d'eux. C'est alors que tous ces cœurs ont trouvé des jouissances indicibles à se communiquer leurs pensées. L'Esprit-Saint avoit allumé dans tous le feu de son amour ; ils se sont rapprochés ; ce contact a augmenté cette divine flamme ; tous ces feux se sont réunis ; ils ont fini par ne plus former tous ensemble qu'un vaste foyer, le foyer d'amour allumé par la main de Dieu même et qu'on a appelé la Congrégation.

Jusqu'ici nous avons suivi l'action de la Providence
% Source : PDF original, vue 038, côté gauche, page manuscrite 76.
disposant tout de loin pour nous préparer dans cette société tant de moyens de salut ; après avoir admiré la bonté de Dieu envers la Congrégation avant qu'elle fût fondée, cherchons dans son histoire de nouveaux motifs pour admirer encore cette bonté divine la soutenant et la conservant comme son enfant chéri.

\begin{center}
Fin du précis historique

\textit{À Deo factum est istud}\footnote{Traduction littérale : « Cela a été fait par Dieu. » Même formule qu'à la p. 46 ; voir Q015. (Note du transcripteur).}
\end{center}
